%% RESUMO NA LINGUA VERNACULA (elemento obrigatorio -- NBR 6028)
%%
%% Regras:
%%   - paragrafo unico, sem recuo, justificado;
%%   - Times New Roman 12, espacamento simples;
%%   - entre 150 e 500 palavras;
%%   - verbo na voz ativa e na terceira pessoa do singular;
%%   - as palavras-chave sao definidas em config/dados.tex.
%%
%% Um bom resumo informa: objetivo geral; principais conceitos e autores da
%% fundamentacao teorica; etapas e acoes da execucao; populacao envolvida e
%% periodo; principais resultados alcancados; conclusoes.

\begin{resumo}
Sistemas de \textit{Retrieval-Augmented Generation} (RAG) dependem da
qualidade da recuperação, e em corpus jurídico-administrativo brasileiro a
busca vetorial pura falha na recuperação de termos exatos e em perguntas que
conectam documentos distintos. Este trabalho compara empiricamente estratégias
de recuperação --- RAG denso, RAG denso-esparso e GraphRAG incremental,
implementado sobre PostgreSQL com pgvector e Apache AGE em armazenamento
único --- e estratégias de geração --- single-pass e verificada multiagente,
no padrão Researcher-Auditor-Adjudicator --- sobre um corpus de documentos
normativos brasileiros --- leis federais, resoluções do Conselho Nacional de
Justiça e instruções normativas do Tribunal de Contas da União. A avaliação
usa um golden dataset de perguntas intra e inter-documentos, as métricas
RAGAS (Faithfulness, Context Precision, Context Recall e Answer Relevancy) e
o teste de Wilcoxon signed-rank pareado com correção de Bonferroni. Na fase
confirmatória, o GraphRAG incremental supera o RAG denso, mas não o RAG
denso-esparso, e a geração multiagente não melhora a Faithfulness nem a taxa
de verificação de citação em relação à single-pass: as hipóteses H1 e H2 não
se confirmam. Uma fase exploratória ampliou a comparação para catorze
condições de recuperação e mostrou que remover o reranking do RAG
denso-esparso eleva a Context Precision sem perda nas demais métricas, e que
a busca esparsa pura (BM25) alcança a maior Context Precision entre as
condições formais. A configuração mais precisa encontrada combina RAG
denso-esparso sem reranking com geração single-pass, mais simples e barata
que a prevista nas hipóteses. O GraphRAG puro, que entrega ao gerador apenas
descrições sintéticas de entidades, ficou abaixo das demais condições, e a
explicação mais sustentada pelos dados é o descasamento entre esse formato e
a métrica de Faithfulness. Entre as limitações do estudo estão o uso de um
único juiz de LLM e uma métrica de recuperação baseada em identificador de
trecho, e não em sobreposição de caracteres.
\end{resumo}
